\begin{abstract}
Dendrites integrate thousands of synaptic inputs, but whether synapses from the same presynaptic neuron are spatially organized on dendrites remains unclear. Using the MICrONS cortical connectome, we analyzed 12 neurons (6 excitatory, 6 inhibitory) receiving 488--2490 synapses each. We first show that synapse spatial structure persists under a soma-distance null model that accounts for known distance-dependent density gradients, confirming higher-order organization beyond distance gradients ($p < 0.015$ for all 12 neurons). We then developed a per-partner clustering test using distance-matched permutation nulls to ask whether synapses from individual presynaptic partners are more spatially compact than expected. Across 496 partners with $\geq 5$ synapses, the median compactness ratio (observed/null mean pairwise geodesic distance) was 0.884, with 87.7\% of partners more compact than expected---a population-wide left shift, not an outlier-driven effect. This enrichment of clustered partners replicated across 10 of 12 neurons (5--11$\times$ enrichment, $p < 0.001$), and was abolished by label-shuffling (real: 24--52\% clustered; shuffled: $< 1\%$), confirming partner-specific spatial organization. The two bipolar cells showed no clustering, consistent with their constrained morphology. Partners with fewer synapses showed stronger compactness but lower detection rates, suggesting distinct spatial strategies for small versus large input groups. These results demonstrate that input-specific dendritic clustering appears to be a common feature of cortical neurons that persists beyond soma-distance gradients, consistent with Hebbian models of spatially organized synaptic connectivity.
\end{abstract}
